\section{Full-source extension with uniformly bounded local multiplicity near an edge}
\label{sec:bounded-boundary-clusters}

The finite boundary divisor in MA-MO-027 can be replaced
by a uniform bound on the total multiplicity in short
rectangles near the edge. Infinitely many boundary zeros,
and infinitely many interior zeros approaching the edge,
are then allowed. The proof keeps each nearby group of
poles in one product; unstable partial-fraction
coefficients are never estimated separately.

Use the actual source settings, forward multipliers
\(A_\omega\), Fourier normalizations, step cutoff and
norms of (MA-MO.700), (MA-MO.705). Thus the setting is
either a nontrivial finite-order Hecke character with
its dual and a structural \(S\), or the rational
principal sector \(S=\{\infty\}\) with its actual
polar completion. Fix \(0<a<1/2\) and \(s\ge0\).
The new divisor hypotheses are
\[
 \begin{gathered}
 |\Re\lambda|\le a
       \quad\hbox{for every completed centered zero of each sector},\\
 \text{there exist }\delta>0,\ \ell>0,\ M\in\mathbb Z_{\ge0}
       \text{ such that for every }T\in\mathbb R,\\
 \sum_{\substack{\lambda:\ a-\delta\le\Re\lambda\le a\\
                         |\Im\lambda-T|\le2\ell}}
                         m_\omega(\lambda)\le M .
 \end{gathered}
 \tag{MA-MO.900}
\]
The bound is imposed for both \(\omega=\eta,\bar\eta\)
in the nontrivial case. Multiplicities are included.
There is no separation condition on distinct zeros.
Shrinking the parameters if necessary, assume
\[
 0<\delta<\tfrac18\min\{a,1/2-a\},
                  \qquad 0<\ell\le1,\qquad
                  0<\epsilon_0<\delta/8 .
 \tag{MA-MO.901}
\]
This shrinking cannot invalidate (MA-MO.900).
All constants concern the fixed field, characters,
place set, \(a,\delta,\ell,M,s\) and cutoff.

\paragraph{MA-MO-030 (bounded local multiplicity suffices for the full matrix).}
Assume (MA-MO.900) and
\[
                         s>2M-1 .
 \tag{MA-MO.902}
\]
When \(M=0\), this inequality holds for every \(s\ge0\).
For every compact smooth test \(g\), the original
\(C_g(P_+-P_-)\) has a unique trace-class extension
on \(H_{\Sigma,a,s}\) in the nontrivial sector and
on \(\mathcal W_{a,s}=H_{\Sigma,a,s}\oplus\mathbb C^2\)
in the rational principal sector. The extensions
are continuous in trace norm on each fixed compact
test-support space. They agree with the original
wide-inverse source formulas.

All completed zero multiplicities are retained.
The evaluated trace and determinant identities
(MA-MO.851)--(MA-MO.858) hold on these actual
operators, including the raw rational determinant's
finite-rank correction and the \(g(0)=0\) condition
for the stated prime formulas.

This is a sufficient condition with a threshold on
the \emph{total local multiplicity} \(M\), not a
claim that the individual boundary-jet condition
alone suffices. The ordinary local zero count
\(O(\log(2+|T|))\) does not imply the fixed bound
in (MA-MO.900). No arithmetic assertion that the
new condition holds is made.

\begin{proof}
\emph{Uniform holomorphic factors after removing nearby zeros.}
Write \(t_n=n\ell\), \(\nu_n=2+|n|\), and let
\(Z_{\omega,n}\) be the multiset of the zeros in
the rectangle appearing in (MA-MO.900) with \(T=t_n\).
Its cardinality with multiplicity \(k_{\omega,n}\)
is at most \(M\). Define
\[
 \begin{aligned}
 P_{\omega,n}(z)&=\prod_{\lambda\in Z_{\omega,n}}(z-\lambda),\\
 b_{\omega,n}(z)&=\frac{P_{\omega,n}(z)}{A_\omega(z)},\\
 \mathcal R_n&=\{z:\ a-\delta/2<\Re z<a+\delta/2,\
                         |\Im z-t_n|<3\ell/2\}.
 \end{aligned}
 \tag{MA-MO.903}
\]
An empty product is one. Every zero of \(A_\omega\)
in \(\mathcal R_n\) is included in \(Z_{\omega,n}\).
The other finite-function zeros, and the finite Euler
zeros on \(\Re z=-1/2\), are outside this rectangle.
The rational pole is outside it by (MA-MO.901).
Thus \(b_{\omega,n}\) is holomorphic in \(\mathcal R_n\).

There are constants \(C_j,D_j\), independent of \(n\),
such that
\[
 \sup_{\substack{a-\delta/4\le\Re z\le a+\delta/4\\
                       |\Im z-t_n|\le5\ell/4}}
                   |b_{\omega,n}^{(j)}(z)|
                         \le C_j\nu_n^{D_j}
                         \qquad(j\ge0).
 \tag{MA-MO.904}
\]
We verify this estimate without assuming separation
inside \(Z_{\omega,n}\). For sufficiently large
\(|t_n|\), use the uncentered finite-\(L\) variable
and a disk centered at \(w_n=2+it_n\), with a
zero-free boundary radius \(12<R_n<13\).
Its total zero multiplicity is \(O(\log\nu_n)\).
All its zeros are nontrivial at these heights.
The entire rectangle \(\mathcal R_n+1/2\)
lies in the disk's inner quarter.

Write \(L_\omega=H_nB_n\), where \(B_n\) is the
finite disk Blaschke product with multiplicities.
The finite-\(L\) polynomial strip bound gives
\(|H_n|\le C\nu_n^D\) on the disk by the maximum
principle. The Euler-product lower bound at
\(2+it_n\), and \(|B_n|\le1\), give
\(|H_n(w_n)|\ge c>0\). Applying Harnack's
inequality to
\(\log(C\nu_n^D)-\log|H_n|\) gives
\[
                        |H_n(w)|^{-1}\le C\nu_n^D
             \quad\hbox{on the inner half disk}.
 \tag{MA-MO.905}
\]
This is the complete finite-disk estimate used
in MT-M18 and (MA-MO.808).

Each normalized Blaschke factor has, up to a phase,
the form
\[
       B_{\rho,n}(w)=
       \frac{R_n(w-\rho)}
        {R_n^2-\overline{(\rho-w_n)}(w-w_n)} .
 \tag{MA-MO.906}
\]
For a zero in \(Z_{\omega,n}\), its numerator
\(w-\rho=z-\lambda\) is canceled exactly by
\(P_{\omega,n}\). The remaining factor
\((w-\rho)/B_{\rho,n}(w)\) is bounded by a
fixed constant on the inner half disk.
There are at most \(M\) such factors.

Every other disk zero is uniformly separated
from \(\mathcal R_n+1/2\). Indeed, either its
centered real part is less than \(a-\delta\),
giving distance at least \(\delta/2\), or
its ordinate differs from \(t_n\) by more
than \(2\ell\), giving distance at least
\(\ell/2\). There are no zeros strictly
to the right of \(a\). Thus the inverses
of all these remaining Blaschke factors
are bounded by a fixed constant depending
only on \(\delta,\ell\). Multiplying
\(O(\log\nu_n)\) such constants costs
only a fixed power of \(\nu_n\).
The deleted Euler factors and their
reciprocals are uniformly bounded on
this closed positive strip. Together
with (MA-MO.905), this proves the
polynomial bound for \(b_{\omega,n}\)
on \(\mathcal R_n\).

For the finitely many bounded-height \(n\)'s,
the canceled quotient is holomorphic on the
same rectangle and bounded on a slightly
smaller fixed rectangle. Its finite bounds
can be incorporated into the constants.
Cauchy estimates on a fixed-radius disk,
using the gap between the two displayed
rectangles, prove (MA-MO.904) for every
derivative. In particular no Laurent
coefficient at an individually isolated
zero has entered the estimate.

\emph{A uniform Sobolev bound for the whole pole product.}
Set \(r=s/2\). Enumerate \(Z_{\omega,n}\),
with repetitions for multiplicity, by
\(\lambda_j=a-d_j+i(t_n+\gamma_j)\),
where \(0\le d_j\le\delta\) and
\(|\gamma_j|\le2\ell\).
For \(k=k_{\omega,n}\ge1\), put
\[
 R_{n,\epsilon}(v)
       =\prod_{j=1}^k
                    (\epsilon+d_j+i(v-\gamma_j))^{-1},
                  \qquad 0<\epsilon\le\epsilon_0 .
 \tag{MA-MO.907}
\]
The character index is suppressed in this local calculation.
For \(k=0\), put \(R_{n,\epsilon}=1\).
For \(k\ge1\), each factor has the elementary
integral representation
\(\int_0^\infty e^{-(\epsilon+d_j)t}
                  e^{i\gamma_jt}e^{-ivt}\,dt\).
Multiplying the Laplace integrals shows
\[
 \begin{gathered}
 R_{n,\epsilon}(v)
       =\int_0^\infty e^{-ivt}e^{-\epsilon t}h_n(t)\,dt,\\
 |h_n(t)|\le\frac{t^{k-1}}{(k-1)!}\quad(t>0).
 \end{gathered}
 \tag{MA-MO.908}
\]
Here \(h_n\) is the convolution of the \(k\)
functions \(\mathbf1_{t>0}e^{-d_jt}e^{i\gamma_jt}\).
The bound is the volume of the \((k-1)\)-simplex,
and is uniform in all \(d_j,\gamma_j\), including
coinciding zeros and zeros on the edge.

The unitary plus Fourier transform of
\(R_{n,\epsilon}\) is
\(\sqrt{2\pi}e^{-\epsilon t}h_n(t)\).
Since \(r>M-\tfrac12\), (MA-MO.908) gives
\[
 \begin{gathered}
 \sup_{\substack{n,\epsilon\\k_{\omega,n}\ge1}}
                    \|R_{n,\epsilon}\|_{H^{-r}}\le C,\\
 R_{n,\epsilon}\longrightarrow R_{n,0}
             \quad\hbox{in }H^{-r}
             \text{ uniformly over the nonempty products}.
 \end{gathered}
 \tag{MA-MO.909}
\]
For the second assertion, the difference has
Fourier transform
\(\sqrt{2\pi}(e^{-\epsilon t}-1)h_n(t)\).
The square of its weighted norm is dominated
by a constant times
\[
 \int_0^\infty
    (1+t^{2M-2})(1+t^2)^{-r}
                         |e^{-\epsilon t}-1|^2\,dt
\]
when \(M\ge1\), with the bounded \(t\le1\)
part interpreted separately. This integrable
majorant and dominated convergence are
independent of \(n\) and of the pole locations.
For \(k=0\) there is no pole product; its
compactly localized contribution below is
handled directly as a smooth function.
Thus \(M=0\) also causes no exception.

\emph{Polynomially controlled localized boundary distributions.}
Choose a smooth partition of unity
\(\sum_{n\in\mathbb Z}\chi_n(u)=1\), with
\(\operatorname{supp}\chi_n\subset
[t_n-\ell,t_n+\ell]\), obtained by translating
and scaling one fixed compact profile.
For \(0<\epsilon\le\epsilon_0\), the exact
localized reciprocal is
\[
 \begin{aligned}
 d_{\omega,n,\epsilon}(u)
   &=\frac{\chi_n(u)}{A_\omega(a+\epsilon+iu)}\\
   &=\chi_n(u)b_{\omega,n}(a+\epsilon+iu)
                                  R_{n,\epsilon}(u-t_n).
 \end{aligned}
 \tag{MA-MO.910}
\]
The smooth compact multiplier
\(\chi_n b_{\omega,n}(a+\epsilon+i\,\cdot)\)
has polynomially bounded derivatives of
every fixed order, uniformly in \(\epsilon\),
by (MA-MO.904). The Sobolev multiplier
estimate (MA-MO.716), applied by duality
to \(H^{-r}\), therefore gives
\[
 \begin{gathered}
 \|d_{\omega,n,\epsilon}\|_{H^{-r}}
                                \le C\nu_n^D,\\
 \|d_{\omega,n,\epsilon}-d_{\omega,n,0}\|_{H^{-r}}
                    \le C\nu_n^D\omega(\epsilon),\qquad
                         \omega(\epsilon)\longrightarrow0 .
 \end{gathered}
 \tag{MA-MO.911}
\]
To see the second bound explicitly, split the
difference between the two smooth multipliers,
which is \(O(\epsilon\nu_n^D)\) in the
required finite derivative norms by Cauchy
estimates, and the pole-product difference
in (MA-MO.909). Both contributions have
the displayed bound after increasing \(D\).
When \(k=0\), the first contribution alone
proves both assertions. All translated
Sobolev norms are the same, so there is
no hidden dependence on the height.

The locally finite sum
\[
                         j_{\omega,+}=\sum_n d_{\omega,n,0}
 \tag{MA-MO.912}
\]
defines a tempered distribution. Indeed,
when paired with a Schwartz function,
insert a fixed translate of a compact
cutoff equal to one on
\(\operatorname{supp}\chi_n\). Its product
with the Schwartz test has \(H^r\) norm
decreasing faster than every power of
\(\nu_n\). This defeats (MA-MO.911)
and proves absolute convergence of the
pairing. The same argument proves
\[
 A_\omega(a+\epsilon+i\,\cdot)^{-1}
       \longrightarrow j_{\omega,+}
                  \quad\hbox{in tempered distributions},\qquad
 A_\omega(a+i\,\cdot)j_{\omega,+}=1 .
 \tag{MA-MO.913}
\]
For the last identity multiply before
taking the limit. The forward multiplier
and its derivatives have uniform polynomial
bounds on these strips, so its product
with a Schwartz test converges in Schwartz
space. Thus (MA-MO.913) is a legitimate
distributional identity.

\emph{Schwartz smoothing sums the entire infinite family in trace norm.}
If \(L(t,u)\in\mathscr S(\mathbb R^2)\), define
\[
       T_{\omega,\epsilon,L}f(t)
            =\sum_n
                 \langle d_{\omega,n,\epsilon}(u),
                                      L(t,u)f(u)\rangle .
 \tag{MA-MO.914}
\]
This defines a trace-class map \(H^r\to H^R\)
for each fixed \(R\ge0\). It converges in
trace norm as \(\epsilon\downarrow0\);
the same conclusion holds if
\(L_\epsilon\to L_0\) in Schwartz space.
Here are the summability details. Insert a
cutoff equal to one on \(\operatorname{supp}\chi_n\)
in the \(u\) variable and translate \(u=t_n+v\).
Translation is isometric on \(H^{\pm r}\).
Every fixed Schwartz seminorm of
\(L(t,t_n+v)\) on the fixed compact \(v\)
support is \(O_N(\nu_n^{-N})\), for every
\(N\), including the necessary polynomial
weights and derivatives in \(t\).
The explicit Hermite-series lemma
(MA-MO.715) and (MA-MO.911) consequently give
\[
  \bigl\|T_{d_{\omega,n,\epsilon},L}\bigr\|_{\mathcal S_1(H^r,H^R)}
                        \le C_{L,N}\nu_n^{D-N}.
 \tag{MA-MO.915}
\]
Choose \(N>D+2\) and sum. The identical
estimate for the differences, using
(MA-MO.911) and the Schwartz convergence
of \(L_\epsilon\), proves the asserted
trace-norm convergence. Only finitely
many Schwartz seminorms are required
for each specified pair of norms.

This constructs the \emph{smoothed} operator.
It does not assert that multiplication by
\(j_{\omega,+}\) is separately defined
as a bounded operator on every \(H^r\) input.

\emph{The boundary value is still the original wide inverse.}
For \(y\in\mathcal E\), let \(F_{\omega,y}\)
be its original series (MA-MO.477) or
(MA-MO.323). Since there are no zeros
strictly to the right of \(a\), the
original wide contour can be shifted
to every \(b>a\); on each such fixed
line it has a positive horizontal gap.
Thus the exact identity (MA-MO.711) holds.

For each Schwartz function \(f\),
(MA-MO.911) also proves that
\(j_{\omega,\epsilon}f\) converges in
\(H^{-r}\), not just in distributions:
the Sobolev multiplier norms of the
localized \(f\)'s decrease rapidly, and
the sum of their product norms is finite.
The same holds for \(u^jf(u)\) for every
integer \(j\ge0\).
Apply this to the uniformly Schwartz
family \(V_by\), with \(b=a+\epsilon\).
Taking Fourier inverses and using local
smooth convergence of \(e^{bx}F_{\omega,y}\)
proves
\[
 \begin{gathered}
      V_aF_{\omega,y}=j_{\omega,+}V_ay,\\
 D_x^j(e^{ax}F_{\omega,y})
             \in L^2((1+x^2)^{-r}\,dx)
                         \qquad(j\ge0).
 \end{gathered}
 \tag{MA-MO.916}
\]
This identifies the boundary prescription
with the original function, without crossing
an unrecorded pole.

For completeness, (MA-MO.916) gives
pointwise polynomial bounds. The elementary
one-dimensional \(H^1\) estimate on
\([x-1,x+1]\), and comparability of
\((1+t^2)^r\) on that interval, imply
\[
 |D_x^j(e^{ax}F_{\omega,y})(x)|
                   \le C_{j,y}(1+x^2)^{r/2}.
 \tag{MA-MO.917}
\]
The original function is smooth, so this
applies to its actual derivatives. At the
right it still has every stronger
exponential weight from the original
wide contour shift. Hence
\(qF_{\omega,y}\in\mathcal E_\omega\),
and is in the actual local source by
the structural-\(S\) inclusion. In the
rational case (MA-MO.714) is again an
actual even additive Schwartz function
vanishing near zero.

\emph{The actual two kernels and their compatible rows.}
Use exactly the removable commutator and
separated crossed kernels (MA-MO.718)--(MA-MO.719).
Their nonsingular factors
\(L_{\omega,b,g}\) and
\(L^{\rm cross}_{\eta,b,g}\)
converge in Schwartz space as \(b\downarrow a\).
The forward multipliers have polynomial
derivative bounds, \(G\) and \(P\) supply
the two rapid directions, and
\(\Re(-z-w)=-2b\le-2a\).
Their proofs do not use a bound for the
reciprocal. The division by \(z-w\)
in the commutator remains removable
before the distribution pairing.
Equation (MA-MO.913) supplies the same
identity \(A j_+=1\).
Applying (MA-MO.914)--(MA-MO.915) gives
trace-class limits
\[
     K_{\omega,a}^{\partial}(g),\qquad
     O_{\eta,a}^{\partial}(g):
               H_{a,s}\longrightarrow H_{a,s},
 \tag{MA-MO.918}
\]
with the appropriate dual input sector
for \(O_{\eta,a}^{\partial}\).
The input reciprocal in that crossed
kernel is \(j_{\bar\eta,+}\), as in
(MA-MO.719), not \(j_{\eta,+}\).

The independent half-line term
\(N_a(g)=C_gM_\phi:H_{\Sigma,a,s}\to H_{a,s}\)
has the trace-class physical kernel
(MA-MO.332). Therefore the first
character row is the actual
\[
       K_{\eta,a}^{\partial}(g)y
              -N_a(g)y+O_{\eta,a}^{\partial}(g)Ry .
 \tag{MA-MO.919}
\]
For every \(y\in\mathcal E\), the original
identity is first valid on every \(b>a\),
with its two fixed original wide inverses
and their actual cutoff source vectors.
The kernel convergence just proved and
\(V_by\to V_ay\) pass that identity
to (MA-MO.919). The actual smooth
theta output is fixed during this
passage. No bounded identity between
whole varying-width Hilbert spaces
is inferred.

The reflected dual row is
\(R\mathscr R_{\bar\eta,a}^{\partial}(g^\vee)R\).
It agrees with the first on the dense
smooth common sector by the actual
Fourier pair and nonpolar Poisson
identity. The paired trace-class
map has range in the closed diagonal
of \(H_{a,s}\oplus H_{-a,s}\), by
density. This proves the claimed
actual nontrivial endomorphism in
the inherited sum norm.

\emph{The rational moments and complete polar matrix.}
The only extra input for the rational
rows is boundedness of the two actual
cutoff moments. Let
\(p_a(x)=\phi(x)e^{(1/2-a)x}\), which
is Schwartz. The functional
\[
 \lambda_0(y)=
       \sum_n\big\langle d_{n,0}(u),
              \widehat p_a(-u)(V_ay)(u)\big\rangle,
                   \qquad\lambda_1(y)=\lambda_0(Ry)
 \tag{MA-MO.920}
\]
is bounded on \(H_{a,s}\), respectively
\(H_{-a,s}\). Indeed the local multiplier
norms of \(\widehat p_a(-u)\) decrease
faster than every power of \(\nu_n\),
and (MA-MO.911) makes this sum of
functional norms absolutely convergent.
For \(y\in\mathcal E\), passage from
the lines \(b>a\) proves precisely
the physical moments (MA-MO.723).
The polynomial left bound
(MA-MO.917) and the stronger right
weights give absolute convergence
at Mellin parameter \(1/2\).
The original zeta quotient has a
zero there, so
\(\mathcal MF_y(1/2)=0\).
Consequently the actual source
\(h_y=u_y+\mathfrak F u_{Ry}\)
has \(h_y(0)=-\lambda_1(y)\) and
\(\int h_y=-\lambda_0(y)\).

All ingredients of the exact
polar-subtracted formula (MA-MO.725)
are now bounded in the required
norms. Its two finite terms remain
\[
 \begin{gathered}
 G(-1/2)\lambda_1(y)e_0+
                   G(1/2)\lambda_0(y)e_1,\\
 \nu_g(y)=
   \bigl(-G(-1/2)\lambda_1(y),
         -G(1/2)\lambda_0(y)\bigr).
 \end{gathered}
 \tag{MA-MO.921}
\]
These signs come from the same
actual theta residue and the two
Gaussian source sections, without
a changed pole subtraction.
The first row is the sum of
(MA-MO.919), in the self-dual case,
and the displayed finite terms.
Its reflected row agrees on the
smooth core. The closed-diagonal
argument gives \(T_g^\partial:H\to H\)
in trace class. Appending the two
fixed genuine source columns gives
exactly the actual full matrix
\[
 \mathscr D_g(y,c)
       =(T_g^\partial y,\nu_g(y))
                             +\sum_{j=0}^1 c_jC_gDs_j .
 \tag{MA-MO.922}
\]
This proves the rational conclusion.

\emph{Continuity, full retention and evaluated traces.}
The original cores are dense, giving
uniqueness. The trace-norm summation
(MA-MO.915) uses a fixed finite
number of Schwartz seminorms for
each target norm. The kernel
seminorms are controlled by finitely
many derivatives of \(g\) on each
fixed compact support. This proves
the asserted trace-norm continuity.

Each right boundary multiplicity \(m\)
is at most \(M\). The completed
functional equation sends a left
boundary zero to a right boundary
zero of the dual sector, with the
same multiplicity; in the rational
case the function is self-dual.
Thus the same bound holds on the
left edge. Consequently (MA-MO.902)
implies \(s>2m-1\). Every boundary
jet is retained by MA-WP; every
interior zero is retained as well.
The actual nontrivial ideal is
invariant and its quotient action
is \(-C_g\), by the same smooth
source formula and density used
in (MA-MO.704). Its restriction
agrees with the unconditional
all-width return of MA-MO-028.
In the rational case the actual
relative operator
\(\mathscr E_g=C_g(1-P_-)-\mathscr D_g\)
has the invariant restriction
\(-L_g\) and quotient \(C_g\).
These are precisely the three
inputs used in the proof of
MA-MO-029: existence of the actual
full matrix, full quotient retention,
and the evaluated positive-source
trace. Its trace, finite-rank
determinant and primitive-prime
calculations therefore give
(MA-MO.851)--(MA-MO.858) here too.
No finite-divisor hypothesis is
used in those algebraic calculations.
\end{proof}

\paragraph{Comparison with the finite-boundary theorem and remaining scope.}
MA-MO-027 is contained in this result.
Choose \(\delta\) smaller than its
interior margin, so only the finitely
many edge zeros remain in the band.
Choose \(\ell\) so small that a
window of length \(4\ell\) contains
at most one distinct edge ordinate,
for either sector. Then \(M\) is
the maximum edge multiplicity, and
(MA-MO.902) is exactly the condition
there. If the edge divisor is empty,
take \(M=0\).

In contrast, (MA-MO.900) permits
infinitely many zeros in the band,
arbitrarily close horizontal
approach to the edge, and arbitrarily
close zero pairs within a window,
provided their \emph{total}
multiplicity stays bounded by \(M\).
The proof's Sobolev exponent pays
for that whole product. It asserts
no necessity of this cluster bound
for a particular arithmetic function.
If the near-edge counts grow with
height, the uniform estimate
(MA-MO.909) at a fixed \(s\) is not
available. The general no-margin
full-retention problem is still open.

The unconditional analytic ingredients
are the polynomial finite-\(L\)
strip bound, Euler-product lower
bound and local \(O(\log|T|)\)
zero count already proved and
located in MT-M18.5--9 and
MA-MO-028. The new assumption is
explicitly the stronger near-edge
count (MA-MO.900), not a conclusion
of that ordinary local count.
The source interfaces and phases
are the already checked
MA-MO-018/022/027/028.

In the inherited cote~1 framework,
scans~079--080 and~111--112 through
BR-TR-01/02, the local multiplicities
remain finite-length divisor data.
The convolution estimate and the
summed actual nuclear operators
supply the additional analytic
structure needed for traces and
determinants. No new arithmetic
positivity or zero-location claim
is inferred from those lengths.
